%%%%%%%%%%%%%%%%%%%%%%%%%%%%%%%%%
%
%       EXERCÍCIO
%
%%%%%%%%%%%%%%%%%%%%%%%%%%%%%%%%%

\ifdefstring{\atividade}{prova}{%
    \renewcommand{\valorquestao}{\ValorQ\ pontos}
}%

\begin{exercicioBanco}[\valorquestao]
Uma indústria automobilística adquire peças de três fornecedores distintos: \(A\), \(B\) e \(C\). O fornecedor \(A\) é responsável por \(20\%\) das peças compradas, o fornecedor \(B\) por \(30\%\) e o fornecedor \(C\) por \(50\%\). Historicamente, a probabilidade de uma peça apresentar algum defeito de fabricação é de \(0{,}05\) para o fornecedor \(A\), \(0{,}03\) para o fornecedor \(B\) e \(0{,}01\) para o fornecedor \(C\).

\begin{enumerate}[(a)]
    \item Qual é a probabilidade de uma peça selecionada aleatoriamente do estoque não apresentar defeito?
    %
    \item Se uma peça escolhida ao acaso apresentou defeito, qual é a probabilidade de ela ter sido produzida pelo fornecedor \(A\)?
\end{enumerate}
%
\end{exercicioBanco}

